\documentclass{article}
\usepackage{amsmath,amssymb,hyperref}
\title{A boundary between completed pure exactness and adjoint admission}
\date{15 November 2026}
\begin{document}\maketitle
\paragraph{Scope.} This is a conditional obstruction to one proposed enlarged doctrine, not to general feedback or directed concurrency. The singleton-module action representation retains native cells and selected synchronization cones and admits right adjoints. The additional requirements below cannot all be imposed on every new arrow.

\paragraph{Assumed parallel doctrine.} Suppose the enlarged process bicategory has a coherent invertibly unital parallel operation $\otimes$, with object unit $I$, and directed interchange
\[
 (f\otimes g);(h\otimes k)\Longrightarrow(f;h)\otimes(g;k).
\]
Assume this comparison is invertible whenever either entire layer consists of admitted pure arrows. Identities are pure. These are the completed-pure-layer exactness conditions used previously; they are stronger than merely having a directed interchange.

\paragraph{Elementary consequence.} For any new endomorphism $F:I\to I$ and admitted pure endomorphism $G:I\to I$, apply interchange to
\[
 (F,1);(1,G)\quad\hbox{and}\quad(1,G);(F,1).
\]
In the first case the second layer is wholly pure; in the second the first layer is wholly pure. Both target pairs are $(F,G)$. The coherent unit constraints therefore give invertible cells
\[
 F;G\ \cong\ F\otimes G\ \cong\ G;F.
\]
No inverse of a general effectful interchange was used. The conclusion follows only from the specified pure-layer inverses and the unit identifications.

\paragraph{Counterexample in the actual proposed completion.} Start with the terminal original pure and process bicategories and identity pure functor. Their local presheaf/enriched-module completion is the ordinary module bicategory, with every completed module admitted pure. At its singleton tensor-unit object, the singleton-action category is Set. The represented pure endomorphism associated with a two-element set acts as
\[
 G(X)=\{0,1\}\times X.
\]
Its right adjoint, admitted by the lax-action enlargement, is $F(X)=X^{\{0,1\}}$. This really is an admitted lax module arrow: its strength sends $(v,(x_0,x_1))$ to $((v,x_0),(v,x_1))$, and arises by the doctrinal mate construction. At the singleton set,
\[
 |F(G(1))|=4,\qquad |G(F(1))|=2.
\]
Hence there is no natural isomorphism $FG\cong GF$, contradicting the preceding necessary consequence. Modifications in the lax-action doctrine have underlying natural transformations, so an invertible modification cannot evade this cardinality obstruction.

\paragraph{Conclusions and nonconclusions.} One cannot extend the old unital parallel operation to all new lax-action arrows while retaining completed-pure-layer exactness against them. The failure already uses the right adjoint of a completed-module pure map, not an arbitrary unrelated functor. It does not contradict the established old-module parallel operation, since $X^2$ is not cocontinuous and is not an old module arrow. It does not invalidate the new cone-preserving representation or its right-adjoint theorem. It does not show that the particular synchronization comparison adjoints needed in a smaller closed process class are incompatible: that requires a separate analysis. Possible repairs must explicitly restrict pure admission or process admission, or weaken the exactness doctrine. Merely changing a coend formula cannot satisfy the inconsistent full set of requirements.

\paragraph{Research consequence.} Keep the general adjoint-plus-cone result as a scoped tool. For the feedback/concurrency unifier, investigate the precise class of comparison adjoints needed and its commutation with the retained pure interface. Any claim of unchanged completed pure exactness must prove this commutation; it cannot be inherited automatically from the old module layer. The obstruction is elementary and no novelty assertion is made.
\end{document}
